% ════════════════════════════════════════════════════════════════════════════
\section{Model}
\label{sec:model}
% ════════════════════════════════════════════════════════════════════════════

\subsection{Principals and capabilities}

\noindent The \textbf{holder} $\Victim$ possesses whatever is needed to reach a
master secret $K$ gating an asset worth $W$; the \textbf{coercer} $\Attacker$
seeks $K$; optional \textbf{delegates} may hold partial capability on $\Victim$'s
behalf. Appendix~\ref{app:notation} tabulates every symbol.

\begin{assumption}[Kerckhoffs, premise~0]\label{ass:kerckhoffs}
$\Attacker$ knows the protocol in full: all algorithms, all source, $\Victim$'s
setup topology, and every in-scope parameter. Only the protocol \emph{secrets}
are hidden. Every result below is proved \emph{against} such an $\Attacker$;
security by obscurity is assumed to contribute nothing.
\end{assumption}

\begin{assumption}[The Wrench--Dolev--Yao (WDY) adversary]\label{ass:power}
The \textbf{Wrench--Dolev--Yao (WDY)} adversary inherits the
Dolev--Yao~\cite{dolev-yao} intruder \emph{on the wire} but departs from it twice. It
is \emph{computational}, not symbolic: perfect cryptography is replaced by the
feasibility line (Stipulation~\ref{ass:feas}), so ``cannot break the cipher''
becomes ``cannot exceed its work budget''. And it is extended \emph{past the wire}
with the physical capabilities Dolev--Yao omits --- device seizure and a bounded
coercion oracle over the holder: the \emph{wrench} the name records. WDY is fixed
by:
\begin{enumerate}[nosep,leftmargin=2.4em,start=0,label=\textbf{D\arabic*.},ref=D\arabic*]
  \item\label{d:kerckhoffs} \textbf{Kerckhoffs.} $\Attacker$ knows the protocol,
        $\Victim$'s setup topology, and every in-scope parameter
        (Assumption~\ref{ass:kerckhoffs}).
  \item\label{d:compute} \textbf{Bounded compute.} $\Attacker$ is bounded by the
        feasibility line (Stipulation~\ref{ass:feas}): $2^{32}$ cheap evaluations
        feasible, plain $2^{64}$ cheap at the border (the softest number),
        $\sftstar$ and $2^{96}$ cheap both infeasible. This is what replaces
        ``unbounded''.
  \item\label{d:frequency} \textbf{Frequency ceiling.} $\Attacker$'s hardware
        evaluates a computation only boundedly faster than $\Victim$'s, so an
        entropy-preserving hash with enough \emph{non-parallelizable} steps
        ($\Hhard$) imposes an arbitrarily long per-evaluation latency on
        $\Attacker$ --- at a proportionally longer one on $\Victim$. This is what
        makes the time-gate meaningful: $\Attacker$ can neither parallelize nor
        out-clock its way past a sequential memory-hard chain.
  \item\label{d:subjugation} \textbf{Complete subjugation.} This is an
        information-security model, not a study of violence against resistance, so
        we do not model $\Attacker$'s coercive power. To keep the argument \emph{a
        fortiori} we assume that from the onset of the attack until release
        $\Attacker$ meets \emph{no resistance whatsoever} and enjoys $\Victim$'s
        \emph{full cooperation}; the sole qualification, \emph{feigned}
        cooperation in the residual window, is of residual importance
        (\S\ref{sec:residual-dr}).
  \item\label{d:seizure} \textbf{Co-located device seizure.} All of $\Victim$'s devices \emph{at
        the attack site} are immediately seized --- vault blobs, time-lock blobs,
        commitment scripts, review schedules. A single WDY encounter has one
        \emph{site}; devices and principals at \emph{other} sites lie beyond its
        reach within the bounded window --- the opening a \emph{geographic}
        defense exploits (\S\ref{sub:fourprops}).
  \item\label{d:explicit} \textbf{Explicit-secret seizure.} All of $\Victim$'s \emph{explicit}
        (non-tacit) authenticating secrets are immediately seized.
  \item\label{d:notacit} \textbf{No tacit seizure.} $\Victim$'s \emph{tacit}
        entropy-carrying secrets --- the points --- are captured \emph{iff} they
        are \emph{available} in transmissible form at the moment of attack
        (materialized, and not wiped before $\Victim$ is subdued): exactly when
        the tacitness fails. When it holds they are not, and by
        \ref{d:subjugation} that is a genuine \emph{incapacity to transmit}, not
        unwillingness --- even a fully cooperating $\Victim$ cannot hand over what
        is recognizable but non-transmissible.
  \item\label{d:noperceptual} \textbf{No perceptual-parameter seizure.} Likewise
        $\Victim$'s perceptual-oracle parameters --- the deep fractals --- are
        seized \emph{iff} available at the moment of attack: they are regenerated
        only by walking the memory-hard chain, so a willing $\Victim$ has nothing
        to dictate. This clause and \ref{d:notacit} are the operational
        restatement, in the adversary model, of the tacit-knowledge premise
        (Assumption~\ref{ass:tkba}).
  \item\label{d:nodelegatecoerce} \textbf{No remote-delegate coercion.} A remote
        delegate --- a party holding rescue capability beyond the attack site
        (\ref{d:seizure}) --- is past $\Attacker$'s \emph{remote} coercive reach:
        from the encounter with $\Victim$, $\Attacker$ cannot compel it to
        surrender or act on its state. Without this a delegated defense is vacuous.
        Real delegates span a spectrum of coercion-compliance; collapsing it to
        ``cannot be reached'' is a deliberate simplification, at the cost of
        reality's shades.
\end{enumerate}
The \emph{only} channel to the tacit material withheld by \ref{d:notacit}--\ref{d:noperceptual} is a
\textbf{bounded coercion oracle}: $\Attacker$ extracts it through the cooperating
$\Victim$ (\ref{d:subjugation}) over a bounded duration $t$ --- the adversary's
\textbf{coercion-time budget}, the maximum wall-clock time for which it can sustain coercion,
calibrated against the incident durations in the registry~\cite{lopp-attacks}.
Cooperation is genuine throughout by \ref{d:subjugation}; the one exception, \emph{feigned}
cooperation, is confined to the residual window (\S\ref{sec:residual-dr}).
\end{assumption}

\subsection{The feasibility line}

Asymptotic ``hard'' and ``easy'' are useless against a fixed, resourced attacker;
these stipulations replace them.

\begin{stipulation}[Feasibility line]\label{ass:feas}
Let a \emph{cheap} evaluation be one call to a non-memory-hard hash $\Hcheap$,
and a \emph{hard} evaluation one call to an hours-scale Argon2-class function
$\Hhard$ whose cost is a large number of \emph{non-parallelizable} steps ---
\emph{memory- and depth}-hard,\footnote{\emph{Depth-hard} is meant more sharply than \emph{sequential} alone: the cost is a non-parallelizable computation with a deliberately large number of steps --- a long sequential \emph{depth}, the parallelism-resistance axis of \emph{depth-robust} memory-hard functions~\cite{depth-robust} --- which $\Attacker$ can neither parallelize away nor out-clock past (\ref{d:frequency}).} so
$\Attacker$'s bounded speed advantage (\ref{d:frequency})
cannot amortize it~\cite{argon2}.
\begin{enumerate}[nosep,leftmargin=2.2em,label=\textbf{S\arabic*},ref=S\arabic*]
  \item\label{s:feasible} \textbf{(feasible).} $2^{32}$ cheap evaluations are within
        $\Attacker$'s budget.
  \item\label{s:infeasible} \textbf{(infeasible).} $2^{64}$ \emph{hard} evaluations --- written
        $\mathbf{64^{*}}$ --- are beyond it.
  \item\label{s:infeasiblec} \textbf{(infeasible).} $2^{96}$ cheap evaluations are beyond it.
\end{enumerate}
The star denotes this memory- and depth-hardness and attaches to \emph{any}
bit-count --- $32^{*}, 64^{*}, 96^{*}, \dots$ --- each meaning that many $\Hhard$
evaluations, not just $64^{*}$. Neither \ref{s:feasible} nor \ref{s:infeasible}
settles the star at $32$ bits, so we state it: \textbf{$32^{*}$ is feasible} ---
memory-hardness alone does not lift a single $32$-bit stage over the floor, which
is why no orbit link commits fewer than $64$ bits (\S\ref{sec:gw}).
\end{stipulation}

\noindent\textbf{The floor.} A work cost is \emph{prohibitive} --- equivalently,
\emph{meets the floor} --- when it satisfies \ref{s:infeasible} \emph{or}
\ref{s:infeasiblec}. The two are independent and we fix \emph{no} ordering between
them: for realistic hardness the per-evaluation $*$ premium is well under $2^{32}$,
so the cheap bound binds, but $*$ is user-tunable and an aggressive setting
reverses that. Every ``floor'' claim names \emph{this} threshold, not a fixed
number, and holds whichever bound binds.

\noindent The star is kept rather than folded into a bit-count because
memory- and depth-hardness has three consequences and a bit-count absorbs only the
first: it inflates the brute-force cost, it imposes a \emph{derivation time} per
evaluation that the coercion window cannot outrun, and it rules out the Grover
speed-up (\S\ref{sec:quantum}). \emph{Entropy} buys $32 \to 64$, but plain
$2^{64}$ \emph{cheap} work still sits at the border of a resourced attacker's
reach, so $64 \to \sftstar$ must be bought by \emph{memory-hardness} --- which is
why the gating function has to be Argon2-class. Plain cheap $2^{64}$ is the
\emph{softest number} throughout: named wherever invoked, never silently treated
as infeasible.

\subsection{The tacit-knowledge premise}

\begin{assumption}[Tacit knowledge, TKBA]\label{ass:tkba}
The secret is a set of fractal locations \emph{recognizable but not
transmissible}: \textbf{TKBA$_1$}, a point is non-transmissible without the
fractal parameters that render it, so $\Victim$ cannot dictate it in the abstract;
and \textbf{TKBA$_2$}, the \emph{deep} fractal parameters are not a hand-overable
artifact, since producing them requires walking a memory-hard chain, before which
later locations do not exist. \S\ref{sec:gw} instantiates the premise.
\end{assumption}

